\documentclass[spanish,12pt,letterpaper,twoside,openright]{report}
\usepackage[T1]{fontenc}
\usepackage[left=3cm, right=3cm, top=2cm, bottom=2cm]{geometry}
\usepackage{graphicx}
\usepackage{mathtools}
\usepackage{amssymb}
\usepackage{amsthm}
\usepackage{xcolor}
\usepackage{babel}
\usepackage{hyperref}
\usepackage{textpos}
\usepackage{xcolor}
\definecolor{UPT}{RGB}{85,18,38}
\begin{document}
	
	%% Logo de la UPT
	\begin{textblock*}{5cm}(-2.5cm,-1cm)
		\includegraphics[width=2cm]{Img/UPT.png}\\
	\end{textblock*}
	
	%% Si tienen Logo de la carrera puedes agregarlo a la carpeta y descomentar las siguientes líneas
	
	%\begin{textblock*}{5cm}(14.5cm,-1cm)
	%	\includegraphics[width=2cm]{UPT.png}\\
	%\end{textblock*}
	
	%%%		Página de Portada		%%%
	\centering
	
	\vspace{1cm}
	
	{\Huge\textbf{Universidad Politécnica de Tulancingo}}
	\textcolor{UPT}{\rule{15.5cm}{1.5pt}}\\
	
	\vspace{2cm}
	
	\textbf{\Large\textcolor{UPT}{Ingeniería en Tecnologías de Manufactura}}\\
	
	\vspace{1cm}
	
	%%		Imagen de Portada		%%
	\begin{figure}[!htb]
		\centering
		\includegraphics[width=\textwidth]{Img/SHM.png}
	\end{figure}
	
	%% 		Nombre de la asignatura		%%
	{\Large	\textbf{Física para Ingeniería}\\}
	
	\vspace{2cm}
	
	%%			Nombre de la práctica		%%
	\textbf{Nombre de la práctica:} COLOCA EL NOMBRE DE LA PRÁCTICA\\
	\vspace{1cm}
	
	%%			Nombre del estudiante		%%
	\textbf{Estudiante:} COLOCA TU NOMBRE AQUÍ
	\vspace{1cm}
	
	%%			Nombre del profesor			%%
	\textbf{Profesor:} Mtro. NOMBRE DEL DOCENTE\\
	\vspace{1cm}
	
	%% 			Fecha actual			%%
	\today
	
	
	
	%%%		Cuerpo de documento		%%%
	%%%		A partir de aquí empieza la siguiente hoja			%%%
	
	\newpage
	
	\textbf{Práctica 1}\\
	\vspace{1cm}
	\raggedright
	\textbf{Nombre de la práctica:} Movimiento Armónico Simple \\
	\textbf{Descripción de la actividad:}  \\
	\textbf{Herramientas necesarias:} Calculadora, computadora, lápiz, borrador  \\
	\vspace{1cm}
	\textbf{Ejercicio(s):} \\
	
	Este es un ejemplo. Por ejemplo, quiero escribir que la ecuación para calcular la fuerza de un resorte debido a su desplazamiento está dada por 
	\begin{align}
		F = -\frac{k}{m} x
	\end{align}
	
	Si desean hacer referencia a la ecuación que escribiste en alguna parte del documento, debes agregarle una etiqueta con el comando \texttt{\textbackslash label\{nombre de la etiqueta\}}, para entenderlo puedes verlo en el siguiente ejemplo, especialmente en el editor de texto. 
	
	\begin{align}
			F = -\frac{k}{m} x \label{ecuacion ley de hooke}
	\end{align}
	
	Con esa etiqueta puedes hacer referencia a la ecuación con el comando \texttt{\textbackslash ref\{ecuacion ley de hooke\}}. Por ejemplo, en la ecuación \ref{ecuacion ley de hooke} podemos ver la ecuación que modela la fuerza de restitución de un resorte con constante de elasticidad $k$ y desplazamiento $x$.\\
	
	
	\vspace{1cm}
	\textbf{Resultados:}  Este es un resultado en negritas $\mathbf{F = 242\ N}$
	
	%% Para resalta un texto puede usar el comanto \textbf{aquí el texto a resalta}
\end{document}